\IfFileExists{stacks-project.cls}{%
\documentclass{stacks-project}
}{%
\documentclass{amsart}
}

% For dealing with references we use the comment environment
\usepackage{verbatim}
\newenvironment{reference}{\comment}{\endcomment}
%\newenvironment{reference}{}{}
\newenvironment{slogan}{\comment}{\endcomment}
\newenvironment{history}{\comment}{\endcomment}

% For commutative diagrams we use Xy-pic
\usepackage[all]{xy}

% We use 2cell for 2-commutative diagrams.
\xyoption{2cell}
\UseAllTwocells

% We use multicol for the list of chapters between chapters
\usepackage{multicol}

% This is generally recommended for better output
\usepackage{lmodern}
\usepackage[T1]{fontenc}

% For cross-file-references
\usepackage{xr-hyper}

% Package for hypertext links:
\usepackage{hyperref}

% For any local file, say "hello.tex" you want to link to please
% use \externaldocument[hello-]{hello}
\externaldocument[introduction-]{introduction}
\externaldocument[conventions-]{conventions}
\externaldocument[sets-]{sets}
\externaldocument[categories-]{categories}
\externaldocument[topology-]{topology}
\externaldocument[sheaves-]{sheaves}
\externaldocument[sites-]{sites}
\externaldocument[stacks-]{stacks}
\externaldocument[fields-]{fields}
\externaldocument[algebra-]{algebra}
\externaldocument[brauer-]{brauer}
\externaldocument[homology-]{homology}
\externaldocument[derived-]{derived}
\externaldocument[simplicial-]{simplicial}
\externaldocument[more-algebra-]{more-algebra}
\externaldocument[smoothing-]{smoothing}
\externaldocument[modules-]{modules}
\externaldocument[sites-modules-]{sites-modules}
\externaldocument[injectives-]{injectives}
\externaldocument[cohomology-]{cohomology}
\externaldocument[sites-cohomology-]{sites-cohomology}
\externaldocument[dga-]{dga}
\externaldocument[dpa-]{dpa}
\externaldocument[sdga-]{sdga}
\externaldocument[hypercovering-]{hypercovering}
\externaldocument[schemes-]{schemes}
\externaldocument[constructions-]{constructions}
\externaldocument[properties-]{properties}
\externaldocument[morphisms-]{morphisms}
\externaldocument[coherent-]{coherent}
\externaldocument[divisors-]{divisors}
\externaldocument[limits-]{limits}
\externaldocument[varieties-]{varieties}
\externaldocument[topologies-]{topologies}
\externaldocument[descent-]{descent}
\externaldocument[perfect-]{perfect}
\externaldocument[more-morphisms-]{more-morphisms}
\externaldocument[flat-]{flat}
\externaldocument[groupoids-]{groupoids}
\externaldocument[more-groupoids-]{more-groupoids}
\externaldocument[etale-]{etale}
\externaldocument[chow-]{chow}
\externaldocument[intersection-]{intersection}
\externaldocument[pic-]{pic}
\externaldocument[weil-]{weil}
\externaldocument[adequate-]{adequate}
\externaldocument[dualizing-]{dualizing}
\externaldocument[duality-]{duality}
\externaldocument[discriminant-]{discriminant}
\externaldocument[derham-]{derham}
\externaldocument[local-cohomology-]{local-cohomology}
\externaldocument[algebraization-]{algebraization}
\externaldocument[curves-]{curves}
\externaldocument[resolve-]{resolve}
\externaldocument[models-]{models}
\externaldocument[functors-]{functors}
\externaldocument[equiv-]{equiv}
\externaldocument[pione-]{pione}
\externaldocument[etale-cohomology-]{etale-cohomology}
\externaldocument[proetale-]{proetale}
\externaldocument[relative-cycles-]{relative-cycles}
\externaldocument[more-etale-]{more-etale}
\externaldocument[trace-]{trace}
\externaldocument[crystalline-]{crystalline}
\externaldocument[spaces-]{spaces}
\externaldocument[spaces-properties-]{spaces-properties}
\externaldocument[spaces-morphisms-]{spaces-morphisms}
\externaldocument[decent-spaces-]{decent-spaces}
\externaldocument[spaces-cohomology-]{spaces-cohomology}
\externaldocument[spaces-limits-]{spaces-limits}
\externaldocument[spaces-divisors-]{spaces-divisors}
\externaldocument[spaces-over-fields-]{spaces-over-fields}
\externaldocument[spaces-topologies-]{spaces-topologies}
\externaldocument[spaces-descent-]{spaces-descent}
\externaldocument[spaces-perfect-]{spaces-perfect}
\externaldocument[spaces-more-morphisms-]{spaces-more-morphisms}
\externaldocument[spaces-flat-]{spaces-flat}
\externaldocument[spaces-groupoids-]{spaces-groupoids}
\externaldocument[spaces-more-groupoids-]{spaces-more-groupoids}
\externaldocument[bootstrap-]{bootstrap}
\externaldocument[spaces-pushouts-]{spaces-pushouts}
\externaldocument[spaces-chow-]{spaces-chow}
\externaldocument[groupoids-quotients-]{groupoids-quotients}
\externaldocument[spaces-more-cohomology-]{spaces-more-cohomology}
\externaldocument[spaces-simplicial-]{spaces-simplicial}
\externaldocument[spaces-duality-]{spaces-duality}
\externaldocument[formal-spaces-]{formal-spaces}
\externaldocument[restricted-]{restricted}
\externaldocument[spaces-resolve-]{spaces-resolve}
\externaldocument[formal-defos-]{formal-defos}
\externaldocument[defos-]{defos}
\externaldocument[cotangent-]{cotangent}
\externaldocument[examples-defos-]{examples-defos}
\externaldocument[algebraic-]{algebraic}
\externaldocument[examples-stacks-]{examples-stacks}
\externaldocument[stacks-sheaves-]{stacks-sheaves}
\externaldocument[criteria-]{criteria}
\externaldocument[artin-]{artin}
\externaldocument[quot-]{quot}
\externaldocument[stacks-properties-]{stacks-properties}
\externaldocument[stacks-morphisms-]{stacks-morphisms}
\externaldocument[stacks-limits-]{stacks-limits}
\externaldocument[stacks-cohomology-]{stacks-cohomology}
\externaldocument[stacks-perfect-]{stacks-perfect}
\externaldocument[stacks-introduction-]{stacks-introduction}
\externaldocument[stacks-more-morphisms-]{stacks-more-morphisms}
\externaldocument[stacks-geometry-]{stacks-geometry}
\externaldocument[moduli-]{moduli}
\externaldocument[moduli-curves-]{moduli-curves}
\externaldocument[examples-]{examples}
\externaldocument[exercises-]{exercises}
\externaldocument[guide-]{guide}
\externaldocument[desirables-]{desirables}
\externaldocument[coding-]{coding}
\externaldocument[obsolete-]{obsolete}
\externaldocument[fdl-]{fdl}
\externaldocument[index-]{index}

% Theorem environments.
%
\theoremstyle{plain}
\newtheorem{theorem}[subsection]{Theorem}
\newtheorem{proposition}[subsection]{Proposition}
\newtheorem{lemma}[subsection]{Lemma}

\theoremstyle{definition}
\newtheorem{definition}[subsection]{Definition}
\newtheorem{example}[subsection]{Example}
\newtheorem{exercise}[subsection]{Exercise}
\newtheorem{situation}[subsection]{Situation}

\theoremstyle{remark}
\newtheorem{remark}[subsection]{Remark}
\newtheorem{remarks}[subsection]{Remarks}

\numberwithin{equation}{subsection}

% Macros
%
\def\lim{\mathop{\mathrm{lim}}\nolimits}
\def\colim{\mathop{\mathrm{colim}}\nolimits}
\def\Spec{\mathop{\mathrm{Spec}}}
\def\Hom{\mathop{\mathrm{Hom}}\nolimits}
\def\Ext{\mathop{\mathrm{Ext}}\nolimits}
\def\SheafHom{\mathop{\mathcal{H}\!\mathit{om}}\nolimits}
\def\SheafExt{\mathop{\mathcal{E}\!\mathit{xt}}\nolimits}
\def\Sch{\mathit{Sch}}
\def\Mor{\mathop{\mathrm{Mor}}\nolimits}
\def\Ob{\mathop{\mathrm{Ob}}\nolimits}
\def\Sh{\mathop{\mathit{Sh}}\nolimits}
\def\NL{\mathop{N\!L}\nolimits}
\def\CH{\mathop{\mathrm{CH}}\nolimits}
\def\proetale{{pro\text{-}\acute{e}tale}}
\def\etale{{\acute{e}tale}}
\def\QCoh{\mathit{QCoh}}
\def\Ker{\mathop{\mathrm{Ker}}}
\def\Im{\mathop{\mathrm{Im}}}
\def\Coker{\mathop{\mathrm{Coker}}}
\def\Coim{\mathop{\mathrm{Coim}}}

% Boxtimes
%
\DeclareMathSymbol{\boxtimes}{\mathbin}{AMSa}{"02}

%
% Macros for moduli stacks/spaces
%
\def\QCohstack{\mathcal{QC}\!\mathit{oh}}
\def\Cohstack{\mathcal{C}\!\mathit{oh}}
\def\Spacesstack{\mathcal{S}\!\mathit{paces}}
\def\Quotfunctor{\mathrm{Quot}}
\def\Hilbfunctor{\mathrm{Hilb}}
\def\Curvesstack{\mathcal{C}\!\mathit{urves}}
\def\Polarizedstack{\mathcal{P}\!\mathit{olarized}}
\def\Complexesstack{\mathcal{C}\!\mathit{omplexes}}
% \Pic is the operator that assigns to X its picard group, usage \Pic(X)
% \Picardstack_{X/B} denotes the Picard stack of X over B
% \Picardfunctor_{X/B} denotes the Picard functor of X over B
\def\Pic{\mathop{\mathrm{Pic}}\nolimits}
\def\Picardstack{\mathcal{P}\!\mathit{ic}}
\def\Picardfunctor{\mathrm{Pic}}
\def\Deformationcategory{\mathcal{D}\!\mathit{ef}}
\usepackage{cleveref}
\crefname{originallemma}{original lemma}{original lemmas}
\Crefname{originallemma}{Original lemma}{Original lemmas}
\makeatletter
\newlabel{lemma-family-associated-closed-specialization}{{0H57}{source}{}{sourcepointer-1}{}}
\newlabel{lemma-family-associated-closed-specialization@cref}{{[originallemma][0][]0H57}{[1][0][]source}{}{}{}}
\makeatother
\crefname{originalequation}{original equation}{original equations}
\Crefname{originalequation}{Original equation}{Original equations}
\makeatletter
\newlabel{equation-cycle-classes}{{0H58}{source}{}{sourcepointer-2}{}}
\newlabel{equation-cycle-classes@cref}{{[originalequation][0][]0H58}{[1][0][]source}{}{}{}}
\makeatother
\crefname{originallemma}{original lemma}{original lemmas}
\Crefname{originallemma}{Original lemma}{Original lemmas}
\makeatletter
\newlabel{more-morphisms-lemma-local-structure-finite-type}{{052E}{source}{}{sourcepointer-3}{}}
\newlabel{more-morphisms-lemma-local-structure-finite-type@cref}{{[originallemma][0][]052E}{[1][0][]source}{}{}{}}
\makeatother
\crefname{originallemma}{original lemma}{original lemmas}
\Crefname{originallemma}{Original lemma}{Original lemmas}
\makeatletter
\newlabel{lemma-relative-cycle-smooth}{{0H5B}{source}{}{sourcepointer-4}{}}
\newlabel{lemma-relative-cycle-smooth@cref}{{[originallemma][0][]0H5B}{[1][0][]source}{}{}{}}
\makeatother
\title{The locus where two relative cycles on the fibres agree is closed}
\author{The Stacks Project Authors}
\date{stacks-project:relative-cycles.tex:local-snapshot-20261012; GNU FDL 1.2 or later}
\begin{document}
\maketitle
\setcounter{subsection}{1}
\noindent\textit{Editorial note.} All delivered mathematical text is the author's own, copied from the pinned Stacks Project chapter file relative-cycles.tex; no statement or proof step was rewritten, reordered or removed. The counted span is the author environment and its own proof environment. References inside the span are handed to the assembler as navigation stubs that carry the canonical tag of the cited tag only; whether each of them is load-bearing is recorded on the stub itself.\par\medskip
\noindent\textit{Allowed standard prerequisites (editorial).}\begin{itemize}
\item Existence of a jointly surjective finite family of proper morphisms putting a given family of cycles on fibres into the image of the cycle-classes construction (relative-cycles chapter, Lemma lemma-family-associated-closed-specialization, canonical tag 0H57). cited at proof lines 1324-1327. Reduces the quasi-compact case to cycles that are in the image of equation-cycle-classes.
\item The cycle-classes construction for families of cycles on fibres (relative-cycles chapter, Equation equation-cycle-classes, canonical tag 0H58). cited at proof line 1327. The image being characterized in the middle reduction.
\item Local structure of morphisms locally of finite type: an étale-local diagram with a fibred square to a scheme Y smooth over the base (more-morphisms chapter, Lemma more-morphisms-lemma-local-structure-finite-type, canonical tag 052E). cited at proof lines 1358-1360. Provides the diagram used to transport the nonvanishing coefficient through flat pullback and proper pushforward.
\item The locus where a relative cycle has nonzero fibre over a smooth morphism is open; more precisely the support behaviour of cycles along a smooth morphism of relative dimension (relative-cycles chapter, Lemma lemma-relative-cycle-smooth, canonical tag 0H5B). cited at proof line 1377. With openness of Y to S, yields the open neighbourhood of s disjoint from the equal locus.
\end{itemize}\medskip
\ifdefined\hypertarget\hypertarget{sourcepointer-1}{}\fi
\noindent\textit{Source reference.} Original lemma 0H57: sources/stacks-chow/tags.txt, source lines 21377--21377. Cited at proof lines 1324-1327: the middle reduction of the proof.\par\medskip
\ifdefined\hypertarget\hypertarget{sourcepointer-2}{}\fi
\noindent\textit{Source reference.} Original equation 0H58: sources/stacks-chow/tags.txt, source lines 21378--21378. Cited at proof line 1327: the image condition in the middle reduction.\par\medskip
\ifdefined\hypertarget\hypertarget{sourcepointer-3}{}\fi
\noindent\textit{Source reference.} Original lemma 052E: sources/stacks-chow/tags.txt, source lines 6567--6567. Cited at proof lines 1358-1360: the diagram transporting the nonzero coefficient.\par\medskip
\ifdefined\hypertarget\hypertarget{sourcepointer-4}{}\fi
\noindent\textit{Source reference.} Original lemma 0H5B: sources/stacks-chow/tags.txt, source lines 21381--21381. Cited at proof line 1377: with openness, gives the neighbourhood disjoint from the equal locus.\par\medskip
\begin{lemma}
\label{lemma-relative-cycles-equal}
Let $f : X \to S$ be a morphism of schemes. Assume $S$ locally Noetherian
and $f$ locally of finite type. Let $r \geq 0$ and
$\alpha, \beta \in z(X/S, r)$. The set $E = \{s \in S : \alpha_s = \beta_s\}$
is closed in $S$.
\end{lemma}

\begin{proof}
The question is local on $S$, thus we may assume $S$ is affine.
Let $X = \bigcup U_i$ be an affine open covering. Let
$E_i = \{s \in S : \alpha_s|_{U_{i, s}} = \beta_s|_{U_{i, s}}\}$.
Then $E = \bigcap E_i$. Hence it suffices to prove the lemma for
$U_i \to S$ and the restriction of $\alpha$ and $\beta$ to $U_i$.
This reduces us to the case discussed in the next paragraph.

\medskip\noindent
Assume $X$ and $S$ are quasi-compact.
Set $\gamma = \alpha - \beta$. Then $E = \{s \in S : \gamma_s = 0\}$.
By Lemma \ref{lemma-family-associated-closed-specialization}
there exists a jointly surjective finite family of proper morphisms
$\{g_i : S_i \to S\}$ such that $g_i^*\gamma$ is in the image
of (\ref{equation-cycle-classes}). Observe that $E_i = g_i^{-1}(E)$ is
the set of point $t \in S_i$ such that $(g_i^*\gamma)_t = 0$.
If $E_i$ is closed for all $i$, then $E = \bigcup g_i(E_i)$
is closed as well. This reduces us to the case discussed in the next
paragraph.

\medskip\noindent
Assume $X$ and $S$ are quasi-compact and $\gamma = \sum n_i[Z_i/X/S]_r$
for a finite number of closed subschemes $Z_i \subset X$
flat and of relative dimension $\leq r$ over $S$.
Set $X' = \bigcup Z_i$ (scheme theoretic union).
Then $i : X' \to X$ is a closed immersion and $X'$
has relative dimension $\leq r$ over $S$. Also
$\gamma = i_*\gamma'$ where $\gamma' = \sum n_i[Z_i/X'/S]_r$.
Since clearly $E = E' = \{s \in S : \gamma'_s = 0\}$
we reduce to the case discussed in the next paragraph.

\medskip\noindent
Assume $X$ has relative dimension $\leq r$ over $S$.
Let $s \in S$, $s \not \in E$. We will show that there exists an
open neighbourhood $V \subset S$ of $s$ such that $E \cap V$ is empty.
The assumption $s \not \in E$ means there exists an integral closed
subscheme $Z \subset X_s$ of dimension $r$ such that the coefficient
$n$ of $[Z]$ in $\gamma_s$ is nonzero. Let $x \in Z$ be the
generic point. Since $\dim(Z) = r$ we see that $x$ is a generic
point of an irreducible component (namely $Z$) of $X_s$.
Thus after replacing $X$ by an open neighbourhood of $x$,
we may assume that $Z$ is the only irreducible component of $X_s$.
In particular, we have $\gamma_s = n[Z]$.

\medskip\noindent
At this point we apply More on Morphisms, Lemma
\ref{more-morphisms-lemma-local-structure-finite-type}
and we obtain a diagram
$$
\xymatrix{
X \ar[dd] & X' \ar[l]^g \ar[d]^\pi & x \ar@{|->}[dd] &
x' \ar@{|->}[l]  \ar@{|->}[d] \\
& Y \ar[d]^h & & y \ar@{|->}[d] \\
S \ar@{=}[r] & S & s & s \ar@{=}[l]
}
$$
with all the properties listed there. Let $\gamma' = g^*\gamma$
be the flat pullback. Note that $E \subset E' = \{s \in S: \gamma'_s = 0\}$
and that $s \not \in E'$ because the coefficient of $Z'$ in $\gamma'_s$
is nonzero, where $Z' \subset X'_s$ is the closure of $x'$.
Similarly, set $\gamma'' = \pi_*\gamma'$. Then we have
$E' \subset E'' = \{s \in S: \gamma''_s = 0\}$ and $s \not \in E''$
because the coefficient of $Z''$ in $\gamma''_s$ is nonzero, where
$Z'' \subset Y_s$ is the closure of $y$. By
Lemma \ref{lemma-relative-cycle-smooth} and openness of $Y \to S$
we see that an open neighbourhood of $s$ is disjoint from $E''$
and the proof is complete.
\end{proof}

\end{document}
